\paragraph{\trMatchingCycle{}。}\label{sub:matching-cycles}
设 $\pi\colon E\to S$ 是一个（一如既往地\trExact{}的）\trLefschetzFibration。取一条\trEmbedding{}路径 $\mu\colon[-1;1]\to\mathrm{int}(S)$，使得 $\mu^{-1}(\mathrm{Critv}(\pi))=\{-1;1\}$。我们可以把它\trSplitting{}成一对具有同一起点的\trVanishingPath{}$\gamma_\pm(t)=\mu(\pm t)$（$t\in[0;1]$），从而得到一对\trVanishingCycle{}$V_{\gamma_\pm}\subset M=E_{\mu(0)}$。让我们先来看一个特别简单的情形，即这两者相等的情形。此时，
\begin{equation}\label{eq:naive-matching-cycle}
\Sigma_\mu=\Delta_{\gamma_-}\cup\Delta_{\gamma_+}
\end{equation}
是\trTotalSpace{}$E$ 的一个\trSmooth\trLagrange\trSubmanifold（由\trLefschetzThimble{}的定义，沿 $\mu$ 的\trParallelTransport{}把诸交集 $\Sigma_\mu\cap\pi^{-1}(t)$ 在一切 $-1<t<1$ 时彼此互映，这在重叠处 $\Delta_{\gamma_-}\cap\Delta_{\gamma_+}=V_{\gamma_\pm}$ 附近给出局部\trChart{} $(-1;1)\times V_{\gamma_\pm}$）。由于是沿边界把两个球\trGluing{}起来的结果，$\Sigma_\mu$ 必然是一个\trHomotopy\trSphere。如果我们假定诸 $V_{\gamma_\pm}$ 的\trFrame{}是\trIsotopy{}的，那么它在\trDifferential{}意义上就是一个标准\trSphere。事实上，只要在这两个\trFrame{}之间选定一个\trIsotopy，就可以把 \eqref{eq:thimble-ball-diffeomorphism} 的那些\trDiffeomorphism{}组合起来，从而得到 $\Sigma_\mu$ 的一个\trFrame。我们将把 \eqref{eq:naive-matching-cycle} 称为\textbf{朴素\trMatchingCycle{}构造}（naive matching cycle construction）。
要求两个\trVanishingCycle{}严格重合，这似乎颇不现实；但人们可以稍稍放松这个假定，代价是对所得\trSphere{}在 $E$ 中的位置失去控制。也就是说，设诸 $V_{\gamma_\pm}$ 作为 $M$ 中带\trFrame{}的\trExact\trLagrange\trSphere{}是\trIsotopy{}的。当情形如此时，我们就说 $\mu$ 是一条\textbf{\trMatchingPath}（matching path）。介于这两个\trVanishingCycle{}之间的任何带\trFrame{}的\trExact\trLagrange\trIsotopy{}$I=(i,j_0,j_1)$，都将类似地称为一条\textbf{\trMatchingIsotopy}（matching isotopy）。初等辛几何告诉我们：存在 $\mathrm{Aut}(M,\partial M)$ 中的一个\trIsotopy{}$(\phi^s)$，使得 $\phi^0=\mathrm{Id}$、$\phi^s\circ i^0=i^s$。由\cref{thm:transport-isotopy-realization}，所给的\trExact\trLefschetzFibration{}存在一个\trDeformation，它\trConstantAlongFibres、并且集中在 $E_{\mu(-1/4)}$ 的一个邻域内，使得沿 $\mu|[-1/2;0]$ 的\trParallelTransport{}映射经此\trDeformation{}变成
\begin{equation*}
h^s_{\mu|[-1/2;0]}=\phi^s\circ h_{\mu|[-1/2;0]};
\end{equation*}
这里，右边的\trParallelTransport{}是就初始的（$s=0$）结构而言的。它对\trVanishingCycle{}的效果是 $V^s_{\gamma_-}=\phi^s(V_{\gamma_-})$，而 $V_{\gamma_+}$ 保持不变。在\trDeformation{}的另一端点 $s=1$ 处，我们有一个\trExact\trLefschetzFibration{}结构，对它而言这两个\trVanishingCycle{}相重合，于是可以经朴素\trMatchingCycle{}构造在\trTotalSpace{}中定义一个\trLagrange\trSphere。由于这两个\trVanishingCycle{}的\trFrame{}也彼此相同（在相关\trVectorSpace{}的认同 $i^0$、$i^1$ 之下），朴素\trMatchingCycle{}便自带一个\trFrame。把\cref{thm:moser-deformation-embedding}应用于反向的\trDeformation，便得到一个从 $(E,\omega^1_E)$ 到 $(E,\omega_E=\omega^0_E)$ 的\trConformallySymplectic{}\trEmbedding；我们用它把朴素\trMatchingCycle{}映回原来的\trSymplecticForm。在带\trFrame{}的\trExact\trLagrange\trSphere{}的\trIsotopy{}意义下，其结果不依赖于 $(\phi^s)$ 的选取、\trLefschetzFibration{}的\trDeformation、以及 $\rho$ 的选取。我们把它称为与 $\mu$ 和 $I$ 相伴的\textbf{\trMatchingCycle}（matching cycle），并记作 $\Sigma_{\mu,I}\subset E$。为了不让记号过于累赘，凡\trIsotopy{}的选取显而易见、或者我们并不在意采取哪一个选取的情形，我们仍将写作 $\Sigma_\mu$。
\begin{remark}\label{thm:matching-isotopy-role-remark}
\trIsotopy{}$I$ 在这一构造中的作用并不容易估量。相对于端点对它作\trDeformation，不影响\trMatchingCycle{}的\trIsotopy{}类。在最低维数 $n=1$ 的情形，\trSurface{}$M$ 中\trExact\trEmbedding{}\trCircle{}的空间具有平凡的 $\pi_1$，因此 $\Sigma_{\mu,I}$ 不依赖于 $I$。高维时相应的问题似乎很难回答。不过我们稍后会看到，由命题 18.21 与命题 18.13 相结合可知，$\Sigma_{\mu,I}$ 作为\trDonaldson--\trFukayaCategory{}中\trObject{}的\trIsomorphism{}类不依赖于 $I$。
\end{remark}
一般而言，并不能保证 $E$ 中给定的任何一个\trLagrange\trSphere{}都能表示成一个\trMatchingCycle（不过，在\trLefschetz{}铅笔的语境下，这类结论的一个\trAsymptotic{}陈述可参见 [16]）。此外，不同的\trMatchingPath{}\trIsotopy{}类可以给出相互\trIsotopy{}的\trMatchingCycle。这不是一种偶见的现象，而是在任何足够复杂的\trLefschetzFibration{}中都可以预期会出现的事情。下面的例子出自\trAuroux：
\begin{example}\label{thm:auroux-matching-example}
为简单起见，考虑最低的维数 $\dim(E)=4$。设在底 $S$ 中有两条\trEmbedding{}路径 $\mu$ 与 $\nu$，连接着两两互异的\trCriticalValue，二者恰好相交一次（显然是在路径内部的某一点处），并且是\trTransverse{}地相交。设 $\mu$ 是一条\trMatchingPath；并且与 $\nu$ 的两半相伴的\trVanishingCycle{}\trTransverse{}相交、且交于单独一点。设 $t_\nu$ 是沿 $\nu$ 的\trHalfTwist。由于沿相交一次的\trCircle{}所作的\trDehnTwist{}满足一个\trBraid{}关系，$\mu'=t_\nu^3(\mu)$ 仍是一条\trMatchingPath，而且结果表明，相伴的\trLagrange{}二维\trSphere{}$\Sigma_{\mu'}$ \trIsotopy{}于 $\Sigma_\mu$（作为一种直接却稍欠直观的办法，人们可以：取一个包含所涉全部\trLefschetzThimble{}的邻域；从\trCellular{}贴附的观点来审视它；发觉它\trSymplectic\trIsomorphic{}于 $T^*S^2$ 的一个\trCompact{}部分；然后应用 [74] 来证明\trLagrange{} $S^2$ 在\trIsotopy{}意义下的唯一性）。这一构造还有一种趣味稍逊的变体：把与 $\nu$ 的两半相伴的\trVanishingCycle{}假定为不交，并取 $\mu'=t_\nu^2(\mu)$。
\end{example}
